% Einheit 1 von 5 – Gleichungen mit zwei Variablen und grafisches Lösen (bank/lineare-gleichungssysteme/e1.jsonl, zusammenbau v0.1)
%% TODO Zweigzeile Teil 1: was hier gelernt wird (Satz in Schülersprache aus dem Einheitstitel) – Einheit 1
\einheitenkopf[e1]{Einheit 1 von 5 · Gleichungen mit zwei Variablen und grafisches Lösen}
\zweigzeile{ab Kl. 9 (am Gymnasium ab Kl. 8) · P10}
\setcounter{aufgabe}{12}

%% TODO \verfahren{…}: Verfahrensüberschrift in Sachsprache fehlt in der Bank (2.3 b)
%% TODO Titel: Ich-kann-Satz fehlt in der Bank (Platzhalter: Kettenname) – Nr. 13
%% TODO Anweisung: ein Satz über der ersten Teilaufgabe fehlt in der Bank – Nr. 13
\begin{aufgabe}{Grafisch}
%% TODO \rechenplatz{n}: Schrittzahl des Grundfalls fehlt in der Bank (nur bei fester Schreibform, 2.3 b)
\begin{teile}
\teil $(2|6)$ – Lösung von welcher Gleichung? I: $x + y = 8$, II: $3x - y = 0$. Kreuze an.\\ \kreuz{nur von I}\\ \kreuz{nur von II}\\ \kreuz{von beiden}\\ \kreuz{von keiner}
\teil $x + y = 5$ – Lösungspaare? Ergänze die Tabelle.

\wertetabelle{x}{y}{0,1,2,3}
\end{teile}
\begin{gleichungsraster}[2]
\gl{\text{$2x + y = 6$ – nach $y$ umstellen.}} &
\end{gleichungsraster}
\begin{teile}
\teil $y - x = 1$ – zeichne die Gerade.

\begin{ksys}[xmin=-5,xmax=5,ymin=-5,ymax=5] \end{ksys}
\teil Schnittpunkt? Zeichne die Geraden zu I: $y = x + 1$ und II: $y = -x + 5$ und lies den Schnittpunkt ab.

\begin{ksys}[xmin=-5,xmax=5,ymin=-5,ymax=5] \end{ksys}
( \leerfeld | \leerfeld )
\teil Probe: Ist $(1|5)$ die Lösung von I: $3x + y = 8$ und II: $x - y = -4$? Rechne die Probe in beiden Gleichungen. \janein
\teil Lösung des Systems? Die Geraden I und II gehören zu den Gleichungen I und II. Lies die Lösung ab.

\begin{ksys}[xmin=-5,xmax=5,ymin=-5,ymax=5,ablesen] \gerade{2}{-3}{I} \gerade{-1}{3}{II} \end{ksys}
( \leerfeld | \leerfeld )
\teil I: $y = 2x + 3$, II: $y = 2x - 1$ – wie viele Lösungen? Kreuze an.\\ \kreuz{genau eine}\\ \kreuz{keine}\\ \kreuz{unendlich viele}
\teil Lösung durch Probieren? I: $x + y = 8$, II: $3x + y = 14$. Setze in I nacheinander $x = 0, 1, 2, \ldots$ ein, trage $y$ ein und prüfe II.

\wertetabelle{x}{y}{0,1,2,3,4,5}
x = \leerfeld , y = \leerfeld
\teil Zelte: Ein Zeltlager hat $14$ Zelte, Zweierzelte und Viererzelte, mit zusammen $44$ Schlafplätzen. $x$ ist die Zahl der Zweierzelte, $y$ die Zahl der Viererzelte. Wie viele Zelte jeder Art sind es? Löse durch Probieren oder grafisch und begründe, dass es keine andere Möglichkeit gibt. \hfill (P10 2016 OS) \\ x = \leerfeld , y = \leerfeld
\end{teile}
\end{aufgabe}

%% TODO \verfahren{…}: Verfahrensüberschrift in Sachsprache fehlt in der Bank (2.3 b)
%% TODO Titel: Ich-kann-Satz fehlt in der Bank (Platzhalter: Kettenname) – Nr. 14
%% TODO Anweisung: ein Satz über der ersten Teilaufgabe fehlt in der Bank – Nr. 14
\begin{aufgabe}{Grafisch, Sek II}
%% TODO \rechenplatz{n}: Schrittzahl des Grundfalls fehlt in der Bank (nur bei fester Schreibform, 2.3 b)
\begin{teile}
\teil Schnittpunkt oder keiner? Kreuze an, wie die gezeichneten Geraden I und II liegen.\\ \kreuz{ein Schnittpunkt}\\ \kreuz{parallel}\\ \kreuz{dieselbe Gerade}

\begin{ksys}[xmin=-5,xmax=5,ymin=-5,ymax=5,ablesen] \gerade{1}{1}{I} \gerade{-2}{4}{II} \end{ksys}
\teil Lösungsgerade? Das System I: $x - y = 2$, II: $3x - 3y = 6$ hat unendlich viele Lösungen. Zeichne alle Lösungen als Gerade und gib die Lösung mit $y = 3$ an.

\begin{ksys}[xmin=-5,xmax=5,ymin=-5,ymax=5] \end{ksys}
x = \leerfeld
\teil Nur diese Lösung? Zeige, dass $x = 1$, $y = 3$ das System I: $2x + y = 5$, II: $x - 3y = -8$ löst, und begründe, dass es keine weitere Lösung gibt. \hfill (Abitur 2023 GK)
\end{teile}
\end{aufgabe}

%% TODO Titel: Ich-kann-Satz fehlt in der Bank (Platzhalter: Kettenname) – Nr. 15
%% TODO Anweisung: ein Satz über der ersten Teilaufgabe fehlt in der Bank – Nr. 15
\begin{aufgabe}{Grafisch – Fehler finden, Begründen}
\begin{teile}
\teil Ben liest bei I: $y = 2x - 1$ und II: $y = -x + 3$ den Schnittpunkt $(1|1)$ ab und macht keine Probe. Finde den Fehler.
\teil Mia stellt um: \rechnung{3x + y &= 5 \\ y &= 3x + 5} Finde den Fehler.
\teil Das System I: $-3x + y = 2$, II: $6x - 2y = -4$ hat unendlich viele Lösungen. Tim zeichnet dafür die Gerade $y = -3x + 2$. Finde den Fehler.
\teil Warum erfüllt der Schnittpunkt zweier Geraden beide Gleichungen des Systems? Begründe.
\teil I: $y = 3x + 1$, II: $y = 3x - 2$ – warum hat das System keine Lösung? Begründe.
\teil Warum haben zwei nicht äquivalente lineare Gleichungen mit zwei Variablen höchstens eine gemeinsame Lösung? Begründe.
\end{teile}
\end{aufgabe}
